\section{Introduction}
\label{sec:introduction}

Anyone who has walked home through a city before sunrise has heard it: a blackbird singing from a lamp post while the sky is still fully dark. The observation is old, and the explanation usually offered is simple. Street lights make the night brighter, so the birds believe that morning has come \citep{hartwell2009urban, okafor2015light}.

That explanation may well be right, but the evidence for it is weaker than its popularity suggests. Bright streets are also loud streets. Traffic noise peaks during the morning rush, and a bird that sings before the rush is heard more clearly than one that sings during it \citep{lindqvist2012noise}. Under this account the birds are not fooled by the light at all. They are avoiding the noise, and light merely happens to be found in the same places.

The two accounts make different predictions, and those predictions can be separated in a city that has quiet bright streets and loud dark ones. This paper reports a three-year study designed around that contrast. We make three contributions:
\begin{itemize}
  \item a dataset of 3{,}912 dawn recordings from 48 sites where illuminance and noise vary almost independently;
  \item an estimate of how far song onset moves per unit of light, with noise, temperature, and date held fixed;
  \item evidence that the effect saturates, which bears directly on how much dimming would be needed to reverse it.
\end{itemize}

We find that light explains most of the shift and noise explains very little. Section~\ref{sec:related} places the study in context, Section~\ref{sec:method} describes the sites and the model, Section~\ref{sec:results} reports the estimates, and Section~\ref{sec:discussion} considers what they mean for lighting policy.
